\section{实验}
\label{sec:evaluation}

本节用于验证论文的主要主张。首先说明评估任务、数据集、指标和基线方法。

\begin{figure*}[!h]
  \centering
  \templatefigure[0.9\linewidth]{定性对比占位图}
  \caption{\textbf{定性对比。} 请替换为有代表性的输出、样例或可视化对比。}
  \label{fig:qualitative_results}
\end{figure*}

\subsection{实验设置}

描述基准协议、基线、指标和统计检验流程。明确哪些结果来自本文方法，哪些是复现自已有工作。

\subsection{主要结果}

报告主要定量结果。数值对比可使用表格，趋势或定性行为可使用图示，并结合引言中的研究主张讨论结果。

\begin{table}[t]
  \centering
  \caption{\textbf{结果示例。} 请将占位数值替换为主要对比结果。}
  \label{tab:results}
  \small
  \setlength{\tabcolsep}{5pt}
  \begin{tabular}{lccc}
    \toprule
    {\sffamily\bfseries\color{hustblue}方法} &
    {\sffamily\bfseries\color{hustblue}准确率} &
    {\sffamily\bfseries\color{hustblue}F1} &
    {\sffamily\bfseries\color{hustblue}提升} \\
    \midrule
    基线 A & 78.4 & 75.9 & -- \\
    基线 B & 80.1 & 77.6 & +1.7 \\
    本文方法 & \hustred{83.6} & \hustred{81.4} & \hustred{+3.5} \\
    \bottomrule
  \end{tabular}
\end{table}

\subsection{消融实验}

分析关键设计选择的作用。每组消融实验都应回答一个明确问题。

\subsection{局限性}

说明已知局限、失败案例及不适用的条件，以便读者正确理解结果。
